%%%%%%%%% PCCP (RSC) JOURNAL TEMPLATE %%%%%%%%%
% Physical Chemistry Chemical Physics (Royal Society of Chemistry), full paper.
% RSC's official LaTeX template is on Overleaf, not on CTAN. This template
% approximates its two-column A4 layout on the standard article class; the
% layout is in pccp-approx.sty (written for this template, not an RSC file).
% RSC says template use is optional and accepts Word, PDF or LaTeX (with a
% PDF) for submission; it retypesets accepted articles (see guidelines.md).
\documentclass[10pt,a4paper,twoside,twocolumn]{article}

%%%%%%%%% PACKAGES - KEEP UNCHANGED %%%%%%%%%
\usepackage{amsmath}
\usepackage{amssymb}
\usepackage{graphicx}
% Tables: booktabs (\toprule, \midrule, \cmidrule, \bottomrule).
\IfFileExists{booktabs.sty}{\usepackage{booktabs}}{}
% Chemistry: mhchem version 4, \ce{...} for formulas and reactions.
\IfFileExists{mhchem.sty}{\usepackage[version=4]{mhchem}}{}
% Quantities: siunitx, \SI/\qty{1.5}{\kilo\joule\per\mole} -> 1.5 kJ mol^-1,
% \si/\unit{\kilo\calorie\per\mole}, \num{1.2e-3}, \SIrange{1}{3}{\angstrom}.
\IfFileExists{siunitx.sty}{\usepackage{siunitx}}{}
% Citations: numbered superscripts, sorted and compressed (RSC style).
\usepackage[super,sort&compress,comma]{natbib}
% RSC-like layout: 8.3 cm columns, framed title/abstract block, sans-serif
% headings, "Fig." captions, page-1 footnotes, table of contents entry. It
% also defines simple fallbacks for \ce, \SI, \si, \num, \qty, \unit and the
% booktabs rules where those packages are missing, and loads hyperref.
\usepackage{pccp-approx}
% Last-resort fallbacks: no effect when mhchem/siunitx or pccp-approx have
% already defined these commands (they normally have).
\providecommand{\ce}[1]{\mbox{#1}}
\providecommand{\num}[2][]{#2}
\providecommand{\si}[2][]{#2}
\providecommand{\SI}[3][]{#2\,#3}
\providecommand{\qty}[3][]{#2\,#3}
\providecommand{\unit}[2][]{#2}

%%%%%%%%% TITLE - PLEASE UPDATE
\title{My Title}

%%%%%%%%% AUTHOR - KEEP THIS PLACEHOLDER EXACTLY AS IT IS
% Printed as "A. Researcher*a" (RSC marks the corresponding author with an
% asterisk); the affiliation and e-mail go to footnote "a" on page 1.
\author{A. Researcher}
\affiliation{Department of Chemistry, Example University, City, Country}
\email{researcher@example.edu}

\begin{document}

%%%%%%%%% TITLE BLOCK %%%%%%%%%
% Prints the framed title, author and abstract block across both columns
% (once the abstract below has been read).
\maketitle

\begin{abstract}
%%%%%%%%% ABSTRACT %%%%%%%%%
Abstract here.
%%%%%%%%% ABSTRACT %%%%%%%%%
\end{abstract}

%%%%%%%%% ESI FOOTNOTE %%%%%%%%%
% Printed on page 1 as the dagger footnote to the title: "Electronic
% supplementary information (ESI) available: <this text> See DOI: ...".
% One sentence naming what the ESI contains; cite it in the text as "ESI".
% (RSC's current author pages and Overleaf template call it "Supplementary
% Information (SI)"; RSC sets the final wording at typesetting.)
\esinote{%
%%%%%%%%% ESI %%%%%%%%%
ESI contents here.
%%%%%%%%% ESI %%%%%%%%%
}

%%%%%%%%% TABLE OF CONTENTS ENTRY (GRAPHICAL ABSTRACT) %%%%%%%%%
% RSC: an image at most 8 cm wide x 4 cm high plus 1-2 sentences (at most
% 250 characters). Printed at the end of the PDF. To fill it, replace both
% occurrences of figures/toc_graphic.png with the most representative figure
% in figures/ (see guidelines.md); keep the placeholder branch.
\begin{tocentry}
\IfFileExists{figures/toc_graphic.png}{%
  \includegraphics[width=8cm,height=4cm,keepaspectratio]{figures/toc_graphic.png}%
}{%
  \fbox{\parbox[c][3.8cm][c]{7.7cm}{\centering
    Graphical abstract placeholder\\[2pt]
    8~cm wide $\times$ 4~cm high}}%
}\par\medskip
%%%%%%%%% TOC TEXT %%%%%%%%%
Table of contents text here.
%%%%%%%%% TOC TEXT %%%%%%%%%
\end{tocentry}

\section{Introduction}
\label{sec:intro}
%%%%%%%%% INTRODUCTION %%%%%%%%%
Introduction here.
%%%%%%%%% INTRODUCTION %%%%%%%%%

\section{Computational details}
\label{sec:methods}
%%%%%%%%% COMPUTATIONAL DETAILS %%%%%%%%%
Computational details here.
%%%%%%%%% COMPUTATIONAL DETAILS %%%%%%%%%

\section{Results and discussion}
\label{sec:results}
%%%%%%%%% RESULTS AND DISCUSSION %%%%%%%%%
Results and discussion here.
%%%%%%%%% RESULTS AND DISCUSSION %%%%%%%%%

\section{Conclusions}
\label{sec:conclusions}
%%%%%%%%% CONCLUSIONS %%%%%%%%%
Conclusions here.
%%%%%%%%% CONCLUSIONS %%%%%%%%%

%%%%%%%%% BACK MATTER (unnumbered, in RSC order) %%%%%%%%%
% State only what the provided materials say; otherwise write
% "To be completed by the authors."
\section*{Author contributions}
%%%%%%%%% AUTHOR CONTRIBUTIONS %%%%%%%%%
Author contributions here.
%%%%%%%%% AUTHOR CONTRIBUTIONS %%%%%%%%%

\section*{Conflicts of interest}
%%%%%%%%% CONFLICTS OF INTEREST %%%%%%%%%
% RSC wording when there are none: "There are no conflicts to declare."
Conflicts of interest here.
%%%%%%%%% CONFLICTS OF INTEREST %%%%%%%%%

\section*{Data availability}
%%%%%%%%% DATA AVAILABILITY %%%%%%%%%
Data availability here.
%%%%%%%%% DATA AVAILABILITY %%%%%%%%%

\section*{Acknowledgements}
%%%%%%%%% ACKNOWLEDGEMENTS %%%%%%%%%
Acknowledgements here.
%%%%%%%%% ACKNOWLEDGEMENTS %%%%%%%%%

%%%%%%%%% REFERENCES %%%%%%%%%
% Printed as "Notes and references" with rsc.bst (copied next to this file).
\bibliographystyle{rsc}
\bibliography{references}
%%%%%%%%% REFERENCES %%%%%%%%%

\end{document}
